\documentclass[12pt]{article}
\usepackage[utf8]{inputenc}
\usepackage[T1]{fontenc}
\usepackage{hyperref}
\title{Thopho Language Resources for K-12 Curriculum Development}
\author{Compiled by Claude Sonnet 5.5 (claude-sonnet-5-5) for Andr\'as Kornai}
\date{}
\begin{document}
\maketitle

\section{Introduction}
Thopho (ISO 639-3 ytp; Glottocode thop1236), also called Black Phula, is a Loloish language of Yunnan. Wikipedia gives about 200 native speakers (2011) in an ethnic population of about 500, in two villages in eastern Guangnan County. Joshua Project reports that Mandarin is increasingly used, especially by younger people, that no writing system is published, and that Bible translation is still needed. Very little material was found, and nothing suggests it is taught in schools.

\section{Primary Educational Resources}
No primers, textbooks or school programmes were found.

\section{Linguistic and Reference Materials}
\begin{itemize}
\item Wikipedia, Thopho language: \url{https://en.wikipedia.org/wiki/Thopho_language}
\item Lund University research portal record that appeared in the search results (content not examined): \url{https://lup.lub.lu.se/record/18677}
\end{itemize}

\section{Religious and Community Resources}
\begin{itemize}
\item Joshua Project, Thopho in China: \url{https://joshuaproject.net/people_groups/20614/CH}
\item Joshua Project language page for ytp: \url{https://joshuaproject.net/index.php/languages/ytp}
\end{itemize}
No community-run language organization was found.

\section{Digital and Audio Resources}
No audio, video, corpus or software resources for Thopho were found.

\section{Curriculum Development Context}
The speaker community is very small and Mandarin is spreading among younger people. There is no published writing system and no teaching material. A curriculum would have to start with documentation and orthography work with the community. The language is not known to be taught in schools.
\end{document}